\documentclass[11pt]{article}
\usepackage[margin=2.3cm]{geometry}
\usepackage{amsmath,amssymb,booktabs}
\allowdisplaybreaks
\input{../build/R06_numbers.tex}
\title{XVI-R R06: loop remainders under the hat-aware rule, $Y$-space completeness,\\
and the \HB{} classes as the invariant evanescent cohomology at $\kappa^3$}
\author{XVI-R (local landmark; imports XVI-Q frozen at Q11; not published)}
\date{}
\begin{document}\maketitle
\begin{abstract}
R05 fixed the probe contraction on the $c^*$ line ($\bar\eta\cdot\bar\eta$, $\hat\eta\cdot\hat\eta$, mixed $=0$), but only
in the $X$-linear part of the ghost rows. (A) We recompute the frozen Q06/Q07 ghost-row \emph{loop remainders} under the
same rule. They \textbf{change}: $ccK$ goes from \LKo{} to \LKn{} terms (\LKd{} differing) and $cc\bar c$ from \LBo{} to
\LBn{} (\LBd{} differing). Every differing term carries exactly one hatted contraction, so the 4D projections are
identical. The drop-$T$ rule is therefore \emph{not} exact for loops (disclosure R06-D1). We propagate the change. The joint
evanescent counterterm fit (full $hhc$ row, both ghost rows, Bose symmetry, ghost equation) stays solvable, with a purely
evanescent particular counterterm verified to annihilate every row. The classification depends only on the homogeneous
system and is unchanged. (B) An independent enumerator shows that the R02--R05 $Y$-basis is exactly the evanescent half of
all local $\kappa^3$, ghost-number $-1$ functionals at the power-counting dimension. We add the \NF{}+\NGgen{}+\NKKc{}
nonvanishing purely-4D generators. Of them, \ZP{} combinations have purely evanescent images. Those images have rank
\ZR{} and all already lie in $T+\text{blind}$, so nothing new becomes trivial. Together with the arguments for total
derivatives, field redefinitions, EOM, Gram identities and parity, this makes $\dim\bar H=\HB$ the \emph{invariant
evanescent cohomology at $O(\kappa^3)$}. ``Physical'' is meant in that precise sense only; no direction is certified as an
observable. Outcome: \textbf{Partial}.
\end{abstract}

\section{Scope and premises}
The premises are XVI-Q Q11 (pinned source-only zip \texttt{1888f327\ldots}) and the R01--R05 layers, whose gates must
ACCEPT. OpenAI.tex and the new sheets are not premises. O.1 stays CLOSED only in its scoped sense; bulk symmetry is not
restored; the insertion is frozen; Route B is not discharged; XV-B and XIV-J are not opened. No frozen code is edited:
the contraction is swapped by a wrapper (\texttt{r05\_blocks.contract\_hat}) for the duration of each call. All R06 code
is in \texttt{r06/}.

\section{(A) The ghost-row loop remainders}
\subsection{Why R05's argument was incomplete}
R05 argued that the drop-$T$ rule is exact for loop expressions because probe dots arise there in
$\eta=\bar\eta+\hat\eta$ pairs. The recomputation below shows that this pairing assumption fails for the loop remainder of the $c^*$-line channel. Single-hatted pieces $T(\cdot,c_i)$ occur there, and the frozen rule annihilates them. We have not isolated which diagram produces them; that is recorded as an observation, not a derivation.

\subsection{Recomputation}
\texttt{r06\_loops.py} evaluates \texttt{ghost\_row\_total(t,0,0)} under both rules:
\begin{center}\begin{tabular}{lcccc}\toprule
row & terms (drop-$T$) & terms (hat-aware) & differing terms & 4D projection\\\midrule
$ccK$ & \LKo & \LKn & \LKd & identical\\
$cc\bar c$ & \LBo & \LBn & \LBd & identical\\\bottomrule\end{tabular}\end{center}
Every difference term has the form $H\,H\,H\,T(\cdot,c_i)$, a single hatted ghost contraction. For example, the $ccK$
difference contains $-\tfrac1{24}(a_K\!\cdot p_1)(c_1\!\cdot p_1)p_1^2\,\hat T(a_K,c_2)$.

\subsection{Propagation to the evanescent fit}
\texttt{r06\_refit.py} re-solves the Q07/Q08 joint system with the hat-aware rule applied to the whole ghost rows. The
unknowns are $\Gamma_{K\!ch}$, $\Gamma_{\bar c ch}$, $\Gamma_{c^*cc}$ and $\Gamma_{hhh}$. The system consists of the
$ccK$, $cc\bar c$ and $hhc$ rows, Bose symmetry, and the ghost equation $X_C=\rho\,T(X_K)$:
\begin{itemize}
\item drop-$T$: solvable; $\dim N$ with the ghost equation $=\QNo$;
\item hat-aware: \textbf{solvable}; $\dim N$ with the ghost equation $=\QNn$. The particular counterterm has
\CTK/\CTC/\CTS/\CTH{} terms in $K/C/S/H$, is purely evanescent, and makes every row vanish identically.
\end{itemize}
The reduced $\bar\Gamma$ homogeneous system reproduces $\dim\bar N_g=\NB$ exactly, as in R05.

\subsection{Effect on earlier released claims}
\begin{itemize}
\item \textbf{Q07 ``all four rows close in 4D'': unaffected.} The 4D ($T\to0$) projections of both remainders are
identical under the two rules, and the 4D counterterms and 4D residual statements involve only these projections.
\item \textbf{Ghost-equation compatibility: unaffected, and re-established.} The hat-aware system including the ghost
equation is solvable with an evanescent particular solution. The frozen Q08 statement used the old remainder; its
conclusion stands, but its witness counterterm is superseded by the R06 one.
\item \textbf{R05 classification: unaffected.} $\bar N_g$, $\bar\Sigma$, $T$ and blind depend only on the $X$-linear
blocks, which R05 already built hat-aware. A change in the inhomogeneous term can only affect existence of a particular
solution, and existence has been verified.
\item \textbf{Frozen certificates} (Q06/Q07 remainder hashes and Q07 witness counterterms) are numerically superseded
for the evanescent part. They are not edited. This is disclosure R06-D1.
\end{itemize}
\textbf{Open:} the $hhc$ loop residual stored in XVI-Q was computed in the full (on-shell $B$) formulation. Paired with
the \emph{reduced} channels ($A$, $D_{EH}$), it gives an unsolvable inhomogeneous system under both rules. The rule change
is therefore not the cause; the residual is formulation-mismatched. A $\bar\Gamma$-formulation loop residual has not
been computed. This does not affect $\bar H$, which is homogeneous.

\section{(B) $Y$-space completeness}
\subsection{The space of local functionals}
At $O(\kappa^3)$, the counterterm-generating functionals $Y$ have ghost number $-1$, three fields and the power-counting
dimension. In the reduced formulation the antighost appears only through $\tilde K=K+\bar c F'(h)$, and $B$ is integrated
out. The candidate field contents are then $\tilde K hh$ (family F), $c^*ch$ (family G) and $\tilde K\tilde Kc$ (family
KKc). No other ghost-number $-1$, three-field content has the right dimension. In momentum space a local functional is a
polynomial in $p_1,p_2$ (with $p_3=-p_1-p_2$) contracted with the polarisation slots. The independent monomials are all
complete pairings of slots with each other or with $p_{1,2}$, with any remaining even momentum degree filled by
$p_i\!\cdot p_j$. Each slot--slot pairing is either barred ($H$) or hatted ($T$).

\textbf{Total derivatives.} On the support of $\delta(p_1+p_2+p_3)$, a total derivative is multiplication by $\sum p_i=0$.
Eliminating $p_3$ removes all such relations, so the pairing monomials in $p_1,p_2$ are a basis modulo total derivatives.

\textbf{Enumeration.} An independent enumerator (\texttt{r06\_complete.enum}) generates all pairings and splits them by
the number of hatted contractions:
\begin{center}\begin{tabular}{lcccc}\toprule
space & all & 4D (no $T$) & evanescent ($\ge1\,T$) & frozen basis\\\midrule
\CountRows\\\bottomrule\end{tabular}\end{center}
In every row the frozen basis coincides exactly with the independent evanescent set. The only gap is the 4D half, which
was missing from $Y$.

\subsection{Enlargement by 4D generators}
For each 4D generator we build $\bar\sigma Y$ with the R03--R05 constructions (F: $D_{EH}$ on all legs, $A_K$,
$\rho T$; G: $(D^T\tilde K)\cdot G$, $\rho T$, $-s_0h\,\delta Y/\delta h$ with $\epsilon_S=-1$; KKc: replace one $\tilde K$ by
$\delta S_{EH}/\delta h$). Each image is split into evanescent coordinates plus a 4D remainder in the 4D basis of each
$X$ sector; nothing is left over. \NF{} F, \NGgen{} G and \NKKc{} KKc generators have nonzero images. The evanescent
parts span rank \EVR{}. The trivial \emph{evanescent} directions are the combinations with vanishing 4D remainder
(the 4D part is fixed separately by the 4D renormalisation). There are \ZP{} such combinations; their images have rank
\ZR{}, lie inside $\bar N_g$ (nilpotency extends), and lie in the already-known $T+\text{blind}$. After enlargement:
$\dim T=\TB$, $\dim(T+\text{blind})=\TPB$, $\dim\bar H=\HB$, with 0 new trivial directions. The R05 basis of 18
representatives transports with rank 18 and none becomes trivial.

\subsection{Field redefinitions and equations of motion}
A field redefinition $h\to h+\kappa^2\Phi(h,c)$ shifts the action by $\frac{\delta S}{\delta h}\Phi$, and in the
antifield formulation by $\bar\sigma(\tilde K\cdot\Phi)$, since $\bar\sigma\tilde K=\delta\bar\Gamma/\delta h$. The same
holds for ghost redefinitions with $c^*$: $\bar\sigma(c^*\cdot\Psi)$. Every such $\tilde K\cdot\Phi$ and $c^*\cdot\Psi$
at this order is a ghost-number $-1$ functional of the F or G type, now included in full. Terms proportional to the EOM
are exactly these images, so ``modulo EOM'' and ``modulo field redefinitions'' are both contained in ``modulo
$\bar\sigma Y$'' on the complete $Y$. Redefinitions that are nonlinear in the antifields contribute through KKc, which is
also complete.

\subsection{Gram identities}
The algebra treats $H$ and $T$ as independent symbols: formal $d$-dimensional contractions with the split
$\eta=\bar\eta+\hat\eta$. In that algebra the pairing monomials are linearly independent by construction. A
four-dimensional Gram identity (antisymmetrisation over five barred indices) is a relation only after restriction to
strictly 4D kinematics. It identifies certain combinations of barred monomials as evanescent, but it does not act on the
formal space in which $\bar\sigma$, $\bar N_g$ and $\bar H$ are defined. Any such identity relates 4D monomials to each
other. Since cohomology is computed on the evanescent quotient with the 4D part fixed separately, a Gram relation can at
most move a direction between the 4D and evanescent labels. No such move occurs among the classes: all 18
representatives have at least one $T$ in every term. The count \HB{} is therefore a statement in the formal
(BMHV-type) algebra. A scheme that imposes Gram identities on the barred subspace would have to repeat the count. This
is the precise sense of ``invariant''.

\subsection{Parity}
$\bar\sigma$, the loop residuals and every ansatz are parity-even (no $\epsilon$ tensor). Parity-odd $Y$ and $X$ form a
separate graded summand that cannot mix with the even one. The 18 even classes are therefore stable under adjoining the
odd sector. The odd evanescent cohomology itself has not been computed.

\subsection{Conclusion of (B)}
On the complete space of local, parity-even, ghost-number $-1$ functionals at $O(\kappa^3)$, the reduced hat-aware
cohomology is $\dim\bar H=\HB$, taken modulo total derivatives, field redefinitions and EOM, and in the formal $H/T$
algebra. Since (A) leaves the homogeneous system untouched and the fit solvable, these \HB{} classes are the
\textbf{invariant evanescent cohomology at this order}. They are ``physical'' only in that precise sense. No class is
certified as an observable direction.

\section{Final counts}
$\dim\bar N_g=\NB$, $\dim T=\TB$, blind $=\BG$, $\dim(T+\text{blind})=\TPB$, $\dim\bar H=\HB$; certified physical
(observable) $=0$.

\section{Disclosures}
\begin{description}
\item[R06-D1] The frozen Q06/Q07 ghost-row loop remainders differ under the hat-aware rule (evanescent only). They are
propagated, the fit is solvable and $\bar H$ is unchanged. Frozen certificates are not edited.
\item[R06-D2] $Y$ is enlarged by the purely-4D generators. There are no new trivial directions.
\item[R06-D3] All R06 code is in \texttt{r06/}; the R01--R05 gates are untouched.
\end{description}

\section{Open items}
\begin{enumerate}
\item The $hhc$ loop residual in the $\bar\Gamma$ formulation (the inhomogeneous reduced fit).
\item The parity-odd evanescent cohomology.
\item Counting in a Gram-reduced scheme, if wanted.
\item p22 ($\kappa^2$ consistency).
\item Bulk symmetry stays not restored; the insertion stays frozen; Route B is not discharged.
\end{enumerate}
\end{document}
